\section{Pregunta 4: Evaluaci\'on de la Inversi\'on de Dependencias en la pr\'actica}

% Responsable: [Integrante 4]

\subsection*{Introducci\'on}

Para juzgar la calidad de dise\~{n}o de la clase OrderProcessor,
es necesario definir el marco normativo que rige la construcci\'on de software orientado 
a objetos y a la estructuraci\'on de arquitecturas limpias.

\textbf{Principio de Inversi\'on de Dependencias (DIP - Dependency Inversion Principle) }

Formulado por Robert C. Martin, el DIP es el \'ultimo de los cinco principios SOLID. Establece que las estructuras de 
software m\'as flexibles y mantenibles son aquellas donde las dependencias del c\'odigo fuente y se refieren exclusivamente a abstracciones y nunca a clases concretas
El principio se desglosa formalmente en 2 reglas :

\begin{itemize}
    \item  Los m\'odulos de alto nivel no deben depender de m\'odulos de bajo nivel, ambos deben de depender de abstracciones.
    \item  Las abstracciones no deben depender de los detalles, los detalles deben de depender de las abstracciones.
\end{itemize}


\textbf{La Regla de Dependencia (The Dependency Rule)}
En la literatura de Clean Architecture, Robert C. Martin extiende el DIP al nivel de componentes
organizados en capas conc\'entricas.
    Las dependencias del c\'odigo fuente solo pueden apuntar hacia adentro, en direcci\'on hacia las pol\'iticas de mayor
nivel y el n\'ucleo de dominio. Nada declarado en un c\'irculo o capa interna (entidades, casos de uso, interfaces) debe conocer, mencionar o importar
elemento alguno perteneciente a un c\'irculo externo.


\subsection*{An\'alisis por estilo arquitect\'onico} 
% -----------------------------------------------------------------
% CASO 1: ARQUITECTURA HEXAGONAL 
% ----------------------------------------------------------------- 
\noindent\textbf{\large Caso 1: Arquitectura Hexagonal (Puertos y Adaptadores - Alistair Cockburn)} 
\subsubsection*{Principios Fundamentales}

    Propuesta por Alistair Cockburn, la Arquitectura Hexagonal nace para evitar que 
los detalles tecnol\'ogicos contaminen la l\'ogica dentral de la aplicaci\'on. Sus pilares esenciales son:
    
\begin{description}
    \item[Distinci\'on Dentro / Fuera (\textit{Inside vs. Outside}):] Elimina la visi\'on apilada tradicional ("arriba y abajo") y establece que solo existe una frontera
relevante: el interior del hex\'agono (la l\'ogica de la aplicaci\'on) contra el exterior del hex\'agono (actores, interfaces de usuario, base de datos y servicios externos).

    \item[Simetr\'ia de actores exteriores:] Desde la perspectiva del n\'ucleo de la aplicaci\'on, una llamada proveniente de una interfaz de usuario, una prueba unitaria automatizada en memoria o una
petici\'on recibida por una red HTTP son sim\'etricas: eventos externos que intentan interactuar con el sistema.

    \item[Puertos (\textit{Ports}):] Un puerto es una interfaz abstracta declarada y poseida por el n\'ucleo. Define un contrato de comunicaci\'on agnostico a la tecnolog\'ia:
    
        \begin{itemize} 
            \item \textbf{Puertos Primarios / Conductores (\textit{Driving}):} Ubicados en la entreda, por donde los actores externos invocan las funcionalidades de la aplicaci\'on.

            \item \textbf{Puertos Secundarios / Conducidos (\textit{Driven})}: Ubicados en la salida, por donde la aplciaci\'on consulta o envia mensajes hacia la infraestructura externa.
        \end{itemize}

    \item [Adaptadores (\textit{Adapters}):] Un adaptador es un componente de infraestructura concreto que rodea a un puerto, actuando como puente de traducci\'on entre el mundo exterior y los contraros del n\'ucleo.

    \item [Ejecuci\'on Aislada  (\textit{Headless Mode}):] El n\'ucleo debe ser capaz de ejecutarse y probarse al 100\% en modo "desacoplado", usando mocks o adaptadores de prueba en lugar de componentes t\'ecnicos reales.

\end{description}

\subsubsection*{An\'alisis de la Violaci\'on en el C\'odigo de OrderProcessor}

En la implementaci\'on descrita en el ejercicio OrderProcessor instancia y llama directamente a la clase AddressValidatorService. Esta instancia
es lo que constituye la dependencia directa y genera el acoplamiento r\'igido con la librer\'ia HTTP subyacente.

\begin{itemize}
    \item \textbf{Invocaci\'on Directa sin Puerto Secundario:} OrderProcessor (el inside del hex\'agono) invoca directamente la clase t\'ecnica AddressValidatorSerice (outside del hex\'agono) sin interponer una interfaz o 
Puerto Secundario en la frontera del n\'ucleo.

    \item \textbf{Infiltraci\'on de Adaptadores de Salida en el N\'ucleo:} El principio de Cockburn exige que el c\'odigo externo no invada el interior del hex\'agono. Al instanciar directamente el adaptador t\'ecnico de salida, las dependencias de la librer\'ia HTTP
invaden la clase del negocio.
    \item \textbf{Ruptura de Modo Desacoplado:}  Al estar soldado a la liber\'ia HTTP real, es imposible ejecutar OrderProcessor en un entorno de pruebas sin contar con conexi\'on a red y un servidor HTTP activo,
destruyendo la testeabilidad de la arquitectura.
\end{itemize}

\noindent\textbf{\large Caso 2: Arquitectura Cebolla/Onion} 
\subsubsection*{Principios Fundamentales}
Propuesta por Jeffrey Palermo, la Aquitectura Cebolla esta diseñada para aplicaci\'ones empresariales de larga
vida, estructurandose en capas conc\'entricas. Sus cuatro principios claves son:

    \begin{itemize}
        \item \textbf{Construcci\'on alrededor de un Modelo de Objetos Independientes:} El centro absoluto de la cebolla es le
modelo de dominio (Domain Model), el cual contiene el estado y comportamento esencial del negocio y solo esta acoplado a si mismo

        \item \textbf{Las Capas Internas Definen las Interfaces, las capas externas la implementan:} La capa central no guarda el c\'odigo de acceso a datos ni
clientes HTTP, \'unicamente contiene los contratos o interfaces. La implementaci\'on concreta de esas interfaces se empuja hacia los bordes exteriores.

        \item \textbf{Direcci\'on de Acoplamiento Estrictamente hacia el Centro:} Todo el c\'odigo de las capas exteriores puede depender de las capas internas, pero el c\'odigo de las capas internas jam\'as puede depender de capas mas externas.
    
        \item \textbf{Compilaci\'on y Ejecuci\'on Independiente del N\'ucleo:} Todo el Application Core debe poder compilarse y ejecutarse de manera completamente aislada de la infraestructura.
    \end{itemize}

\subsubsection*{An\'alisis de la Violaci\'on en el C\'odigo de \texttt{OrderProcessor}}
    \begin{itemize}
        \item \textbf{Contaminaci\'on del Centro por la Infraestrutura:} En lugar de empujar la tecnolog\'ia HTTP hacia el borde exterior del a cebolla, la implementaci\'on de Order Processor trae la llamada HTTP al coraz\'on de la aplicaci\'on.
        
        \item \textbf{Inversi\'on Negativa del Acoplamiento:} El acoplamiento viaja desde el centro (Application Core) hacia afuera( infrastructure),
        violando directamente la norma de que "toda acoplamiento es hacia el centro".

        \item \textbf{Acoplamiento a Tecnolog\'ias Cambiantes:} Palermo enfatiza que la infraestructura es lo que cambia con mayor frecuencia en la industria. Al tener a OrderProcessor acoplado a
        AddressValidatorService y a su liber\'ia HTTP, cualquier actualizaci\'on o descontinuaci\'on de la librer\'ia obligar\'a a modificar el c\'odigo fuente de la l\'ogica de procesamiento de pedidos.
    \end{itemize}

\noindent\textbf{\large Caso 3: Arquitectura Limpia/Clean Architecture} 
\subsubsection*{Principios Fundamentales}
Robert C. Martin sintetiza las ideas de Hexagonal y Onion definiendo cuatro c\'irculos conc\'entricos principales

\begin{itemize}
    \item \textbf{Entidades  (\textit{Entities}):} Reglas de negocio cr\'iticas y globales
    \item \textbf{Casos de Uso  (\textit{Use Cases/ Interactors}):} Pol\'iticas espec\'ificas de la aplicaci\'on que coordinan el flujo de datos
    \item \textbf{Adapatadores de Interfaz  (\textit{Interface Adapters}):} Presentadores, controladores y gateways que convierten datos.
    \item \textbf{Frameworks y Drivers:} Herramientas externas, bases de datos, UI y clientes HTTP.
\end{itemize}

Sus principios rectores son: 
\begin{itemize}
    \item \textbf{N\'ucleo como la Pol\'itica de Mayor Nivel:} Las reglas de negocio ocupan la posici\'on mas privilegiada de la jerarqu\'ia y deben de ser protegidas de cualquier cambio en las capas externas
    \item \textbf{La Infraestructura como un Plugin:} Las bases de datos, la interfaz web o las librer\'ias HTTP deben comportarse como meros complementos intercambiables (plugins) que se conectan al sistema sin alterar el n\'ucleo.
    \item \textbf{Cruce de Fronteras mediante Polimorfismo:}  Para que el flujo de control ejecute c\'odigo en la infreaestructura sin violar la Regla de Dependencias, see utiliza el DIP: el Caso de Uso llama a un puerto/interfaz interno, y la infraestructura lo implementa.
\end{itemize}

\subsubsection*{An\'alisis de la Violaci\'on en el C\'odigo de OrderProcessor}
\begin{itemize}
    \item \textbf{Violaci\'on de la Direcci\'on de la Regla de Dependencia:} OrderProcessor reside en la capa de Casos de Uso (pol\'itica de alto nivel), AddressValidatorService reside junto a la librer\'ia HTTP en la capa de Frameworks y Drivers
    (mecanismos de bajo nivel). La flecha de dependencia apunta directamente hacia afuera, rompiendo el principio arquitect\'onico fundamental.
    \item \textbf{Destrucci\'on del Concepto de Plugin:} Debido a que la clase OrderProcessor utiliza la instanciaci\'on directa, el servicio de validaci\'on de direcci\'ones no es un plugin reemplazable, sino una pieza "soldada" dentro del m\'odulo central de la aplicaci\'on.
\end{itemize}
\subsubsection*{¿Qu\'e direcci\'on deber\'ian de tener las dependencias del c\'odigo fuente?:}
    Las dependencias del c\'odigo fuente deben apuntar estrictamente de afuera hacia dentro. El c\'odigo de alto nivel (dominio/casos de uso) jam\'as debe de mencionar, importar o instancias clases pertenecientes a m\'odulos de bajo nivel (bases de datos, clientes HTTP, UI)

    
\subsection*{¿Por qu\'e la implementaci\'on actual no cumple?: }
    \begin{itemize}
        \item   Porque la flecha de dependencia apunta en \textbf{direcci\'on contraria (de adentro hacia afuera)}.
        \item La clase de alto nivel \texttt{OrderProcessor} (n\'ucleo de negocio) conoce, depende e instancia directamente a la clase t\'ecnica de bajo nivel \texttt{AddressValidatorService}
    (infraestructura), quedando atada a su librer\'ia HTTP subyacente.

    \end{itemize}


\subsection*{ Redise\~{n}o Arquitect\'onico e Inversi\'on de Dependencias}
Se aplica el Principio de Inversi\'on de Dependencias (DIP) junto con el patr\'on de interfaz separada.

\begin{lstlisting}[language=CSharp, caption={Ejemplo de c\'odigo}]
// Puerto / Interfaz declarada en el centro (propiedad del negocio)
public interface IAddressValidator
{
    bool ValidateAddress(Address address);
}

//Caso de Uso / l\'ogica de Negocio que depende \'unicamente de la abstraccion
public class OrderProcessor
{
    private readonly IAddressValidator _addressValidator;

    // Inyecci\'on de Dependencias por Constructor
    public OrderProcessor(IAddressValidator addressValidator)
    {
        _addressValidator = addressValidator
        ?? throw new ArgumentNullException(nameof(addressValidator));
    }
    public void ProcessOrder (Order order)
    {
        // La regla de negocio invoca el contrato sin conocer detalles de red
        bool isValid = _addressValidator.ValidateAddress(order.ShippingAddress);

        if(!isValid)
        {
            throw new InvalidOperationException("La direcci\'on de envio no es valida");
        }
        // l\'ogica adicional para procesar el pedido.....
    }
}
\end{lstlisting}

\begin{lstlisting}[language=CSharp, caption={Ejemplo de c\'odigo}]
//Adaptador concreto en la infraestructura que implementa la interfaz del n\'ucleo

public class AddressValidatorService: IAddressValidator
{
    private readonly HTTPClientLibary _httpClient;

    public AddressValidatorService(HttpClientLibrary httpClient)
    {
        _httpClient = httpClient;
    }

    public bool ValidateAddress(Address addres)
    {
        // Detalle t\'ecnico: traducci\'on de la llamada y comunicaci\'on HTTP
        var response = _httpClient.Post("/api/address/validate", address);
        return reponse.IsSuccess;
    }
}
\end{lstlisting}

\subsection*{An\'alisis por casos y ubicaci\'on de componentes}



\noindent\textbf{\large Caso 1: Arquitectura Hexagonal}

\begin{itemize}
    \item \textbf{N\'ucleo de la Aplicaci\'on (Interior del Hex\'agono /Inside):}
    \begin{itemize}
            \item IAddressValidator: Reside dentro del hex\'agono actuando como un Puerto Secundario / Conducido (Driven Port). Define la conversaci\'on abstracta necesaria sobre la frontera del negocio.

            \item OrderProcessor: Reside dentro del hex\'agono como la l\'ogica central de la aplicaci\'on que utiliza el puerto secundario
    \end{itemize}

        \item \textbf{Capa de Infraestructura (Exterior del Hex\'agono / Outside):}
    \begin{itemize}
            \item AddressValidatorSerice: Reside en el exterior del hex\'agono como un adaptador Secundario / Conducido (Driven Adapter). Implementa el puerto e interactua directamente con el dipositivo o servicio HTTP externo
    \end{itemize}
\end{itemize}

\noindent\textbf{\large Caso 2: Arquitectua Cebolla / Onion}
\begin{itemize}
    \item \textbf{N\'ucleo de la Aplicaci\'on (Application Core / Anillos Internos)}: 
    \begin{itemize}
        \item IAddressValidator: Reside en las capas internas el n\'ucleo de aplicaci\'on (Application Core) Palermo establece expresamente que las
        "las capas internas definen las interfaces"

        \item OrderProcessor: Reside en los servicios del n\'ucleo de aplicaci\'on,
        acoplado \'unicamente a la interfaz y al modelo de objetos independientes
    \end{itemize}

    \item \textbf{Capa de Infraestructura (Anillo Exterior / Infraestructure):}:
    \begin{itemize}
        \item AddressValidatorService: Reside en el borde mas externo de infraestructura. Implementa la interfaz del n\'ucleo e integra la librer\'ia HTTP
        concreta.

        \item En tiempo de ejecuci\'on, un contenedor de Inversi\'on de Control (IoC) resuelve la clase t\'ecnica del a infraestructura e inyecta su instancia en 
        el constructor de OrderProcessor
    \end{itemize}
\end{itemize}

\noindent\textbf{\large Caso 3: Arquitectua Limpia/ Clean Architecture}
\begin{itemize}
    \item \textbf{N\'ucleo de Dominio y Casos de Uso (C\'irculos conc\'entricos internos):}
    \begin{itemize}
        \item IAddressValidator: Reside en la capa de Casos de Uso actuando como una interfaz de l\'imite de salida.
        \item OrderProcessor: Reside en la capa de Casos de Uso, coordinando el flujo de datos y aplicando las reglas de negocio.
    \end{itemize}

    \item \textbf{Capa de Infraestructura y Drivers (C\'irculos Conc\'entricos Externos):}
    \begin{itemize}
        \item AddressValidatorService: Reside en el c\'irculo exterior de Adaptadores de Interfaz / Frameworks and Drivers. Implementa el puerto
        de salida del caso de uso mediante polimorfismo din\'amico.
    \end{itemize}
\end{itemize}

La restructuraci\'on realizada se justifica mediante la segunda regla del Principio de Inversi\'on de Dependencias (DIP), la cual establece que las abstracciones 
no deben depender de los detalles , sino que los detalles deben depender de las abstracciones. 

\begin{itemize}
    \item \textbf{Inversi\'on de la Flecha de Compilaci\'on:} En la implementaci\'on original defectuosa, OrderProcessor (negocio) conoc\'ia y depend\'ia de AddressValidatorService
    (detalle t\'ecnico). En el redise\~{n}o, la clase t\'ecnica AddressValidatorService debe de importar e implementar la interfaz IAddressValidator. Por lo cual el detalle t\'ecnico de la infraestructura ahora
    depende explicitamente de la abstracci\'on definida por el negocio.

    \item \textbf{Aislamiento de la Volatilidad de los Detalles:} Los clientes o librer\'ias HTTP son detalles vol\'atiles que cambian frecuentemente por actualizaci\'ones tecnol\'ogicas. 
    Al hacer que la infraestructura depende de la interfaz abstracta IAddressValidator, los cambios t\'ecnicos en la llamada HTTP quedan totalmente contenidos en la capa eterna y no provocan ning\'un impacto ni recompilaci\'on en el n\'ucleo de la aplicaci\'on.

    \item \textbf{Inyecci\'on de Dependencias por Constructor (DI):} La clase OrderProcessor ya no utiliza instanciaci\'on directa. Mediante la Inyecci\'on de Dependencias por contructor, el negocio declara que abstracci\'on necesita para operar, delegando la provisi\'on de la implementaci\'on concreta a un contenedor IoC en el incio del a aplicaci\'on.
\end{itemize}


\subsection*{Pruebas Aisladas y Mantenibilidad del Software}
\noindent\textbf{\large Impacto en las Pruebas Aisladas}
El redise\~{n}o mediante la Inversi\'on de Dependencias (DIP) y la Inyecci\'on de Dependenciaas por Constructor (DI) resuelve de manera definitiva
los problemas de testeabilidad de la clase OrderProcessor:

\begin{itemize}
    \item \textbf{Eliminaci\'on de la Dependencia de Red y Servidores Externos}
    \begin{itemize}
        \item \textbf{Dise\~{n}o Original:} Para probar la l\'ogica de OrderProcessor, el entono de pruebas estaba obligado a realizar peticiones 
        HTTP reales a un servidor externo. Si el servidor ca\'ia, la red fallaba o la API externa cambiaba, las pruebas unitarias fallaban aunque el c\'odigo de negocio estuviera correcto
        \item \textbf{Dise\~{n}o Rediseñado:} Al depender de la interfaz IAddressValidator, OrderProcessor se vuelve completamente independiente al mecanismo de validaci\'on real y al a red.
    \end{itemize}
\end{itemize}

\subsection*{Uso de Objetos Simulados (Mocks/ Stubs)}
En el arn\'es de pruebas unitarias, se crea un objeto Mock que implementa la interfaz IAddressValidator y se inyecta directamente en el constructor de OrderProcessor. Esto permite simular escenarios de prueba en memoria sin realizar llamadas de red:

\begin{lstlisting}[language=CSharp, caption={Ejemplo de c\'odigo}]
[Fact]
public void ProcessOrder_DireccionInvalida()
{
    var mockValidator = new Mock<IAddressValidator>();

    mockValidator.Setup(v => .ValidateAddress(It.IsAny<Address>())).Returns(false);

    var processor = new OrderProcess(mockValidator.Object);
    var order = new Order();

    Assert.Throws<InvalidOperationException>(() => processor.ProcessorOrder(order));
}

\end{lstlisting}
\subsection*{Cumplimiento de las Propiedades de Pruebas Unitarias de Calidad (FIRST)}
\begin{itemize}
    \item \textbf{Velocidad (Fast):} Las pruebas se ejecutan en milisegundos en memoria, al no depender del tiempo de respuesta I/O de la red.
    \item \textbf{Aislamento (Isolated):} Se prueba exclusivamente la l\'ogica de procesamiento del pedido sin interferencia de componentes de infraestructura.
    \item \textbf{Repetibilidad (Repeatable):} Las pruebas producen exactamente el mismo resultado en cualquier entorno (m\'aquina de desarrollador, servidor de integraci\'on continua CI/CD) sin requerir conexi\'on a Internet
\end{itemize}

\subsection*{Impacto en la Mantenibilidad del Software}
La mantenibilidad mide la facilidad con la que el software puede modificarse para corregir errores, adaptarse a nuevos requisitos o evolucionar tecnol\'ogicamente. El redise\~{n}o optimiza este atributo de calidad a traves de dos principios fundamentales de SOLID:

\begin{itemize}
    \item \textbf{Principio de \'Unica responsabilidad (SRP):}
    \begin{itemize}
        \item OrderProcessor adquiere una sola raz\'on para cambiar: modificaciones en las reglas de negocio relativas al procesamiento de pedidos.
        \item Si el servicio externo de validaci\'on de direcci\'ones cambia en sus endpoints, actualiza su formato JSON o migra de librer\'ia HTTP, la clase OrderProcessor no sufre ning\'un cambio ni requiere recompilaci\'on. El impacto
        del cambio queda 100\% contenido dentro de la clase AddressValidatorService en la capa de infraestructura.
    \end{itemize}

    \item \textbf{Principio Abierto/Cerrado (OCP)}
    El sistema queda abierto a la extensi\'on pero cerrado a la modificaci\'on:

    \begin{itemize}
        \item Si en el futuro la empresa decide cambiar el proveedor de validaci\'on de direcci\'ones o utilizar una base de datos local para la validaci\'on, no se modifica una sola linea de c\'odigo dentro de OrderProcessor.
        \item Basta con crear una nueva clase en la infraestructura que implemente la interfaz IAddressValidator e inyectarla mediante el contenedor de Inversi\'on de Control (IoC) 
    \end{itemize}

    \item \textbf{Reducci\'on de la Deuda T\'ecnica}
    Como enfatizan los autores de las arquitecturas Hexagonal, Cebolla y Limpia, proteger el n\'ucleo de aplicaci\'on de la volatidad de la infraestructura reduce drasticamente los costos de mantenimiento y extiende el tiempo de vida \'util del sistema.
\end{itemize}